% Modernised reading — fonds Alexandre Grothendieck, shelfmark 161-1, batch 1
% (pages 1-19).
%
% This is an INTERPRETATION, not a transcription. It re-reads the pages in
% today's notation and vocabulary, states the hypotheses the manuscript leaves
% implicit, and drops the critical apparatus so that the mathematics reads
% straight through. Everything the brackets and underlines carried — a doubtful
% reading, a notation he used differently, a missing passage — moves into a
% footnote.
%
% It is held to being mathematically correct as it stands, not merely faithful
% to the page. Where the manuscript is loose or mistaken, the true statement is
% what appears here, and the footnote says what the page has. In any dispute of
% substance the transcription governs: whoever cites Grothendieck must cite
% batch-01.fr.tex.
%
% Pass: Fable 5 (claude-fable-5), 2026-08-08 — first pass, unchecked against the
% pages by a human.
\documentclass[11pt,a4paper]{article}
% Le préambule commun à toutes les transcriptions du fonds Grothendieck.
%
% Il tient en une idée : **l'incertitude se compose, elle ne se résout pas.**
% Une transcription qui lisse un mot douteux a détruit la seule chose pour
% laquelle on la faisait — un lecteur doit pouvoir savoir, à chaque endroit,
% ce qui était sur le feuillet et ce qui a été ajouté. Les six macros qui
% suivent sont donc l'appareil critique, et non de la décoration.
%
% Les mêmes commandes sont reconnues par `scripts/render.mjs`, qui produit la
% vue de lecture du site. Une macro ajoutée ici doit l'être aussi là-bas,
% faute de quoi le rendu échouera bruyamment — ce qui est voulu : un rendu
% silencieusement faux est pire qu'un rendu absent.

\ProvidesPackage{grothendieck}[2026/08/08 Transcriptions du fonds Grothendieck]

\usepackage{amsmath,amssymb,amsthm}
\usepackage{mathrsfs}
\usepackage{stmaryrd}
% Les diagrammes commutatifs sont partout dans ce fonds : c'est la langue
% même de la géométrie algébrique de Grothendieck, et une transcription qui
% les rendrait en prose ne transcrirait rien.
\usepackage{tikz-cd}
% Français par défaut — c'est la langue des feuillets — mais l'anglais est
% chargé aussi : la traduction et le résumé sont des documents anglais, et la
% typographie française (espaces avant les deux-points) y serait une faute.
% Ces documents basculent par \selectlanguage{english} en tête de corps.
\usepackage[english,french]{babel}
\usepackage{fontspec}
\usepackage{geometry}
\geometry{a4paper,margin=25mm,marginparwidth=28mm}
\usepackage{marginnote}
\usepackage{xcolor}
% ulem, not soul. `soul`'s \st and \ul are built on a character-by-character
% scan that XeLaTeX's font machinery breaks: the document fails with an
% undefined control sequence rather than degrading. ulem does the same two jobs
% and is Unicode-engine safe. `normalem` stops it hijacking \emph, which the
% transcriptions use for Grothendieck's own emphasis.
\usepackage[normalem]{ulem}
\usepackage{hyperref}
\hypersetup{colorlinks=true,linkcolor=black,urlcolor=black,pdfborder={0 0 0}}

% --- Métadonnées du lot -----------------------------------------------------
% Reprises telles quelles par le script de rendu, qui les affiche en tête de
% la vue de lecture. Elles fixent ce que le fichier prétend couvrir : sans
% elles, une transcription flotte sans point d'attache dans un fonds de seize
% mille feuillets.

\newcommand{\folder}[1]{\gdef\grothFolder{#1}}
\newcommand{\batch}[1]{\gdef\grothBatch{#1}}
\newcommand{\leaves}[2]{\gdef\grothFirst{#1}\gdef\grothLast{#2}}
\newcommand{\foldertitle}[1]{\gdef\grothFoldertitle{#1}}
\gdef\grothFolder{?}\gdef\grothBatch{?}\gdef\grothFirst{?}\gdef\grothLast{?}\gdef\grothFoldertitle{}

% --- Le résumé --------------------------------------------------------------
%
% Il ouvre la lecture modernisée, sous le titre « Résumé », et il est composé
% en petit. Ce n'est pas un ornement : c'est le seul passage du document qui
% ne soit pas des mathématiques, et un lecteur doit pouvoir le reconnaître
% avant de l'avoir lu — puis le sauter s'il connaît déjà le sujet, ou s'y
% arrêter s'il ne le connaît pas. Le filet qui le suit marque la reprise du
% corps.

\newenvironment{resume}
  {\small\setlength{\parskip}{0.45em}}
  {\par\medskip\hrule\medskip}

% --- Mots-clés ---------------------------------------------------------------
%
% En anglais, en clôture du résumé de la lecture modernisée : le vocabulaire
% moderne sous lequel chercher ce que les feuillets construisent. Le manifeste
% du site (`npm run manifest`) les extrait pour étiqueter la cote entière —
% cette ligne est la source unique des tags, il n'y a pas de fichier à côté.
\newcommand{\keywords}[1]{\par\smallskip\noindent
  {\footnotesize\textsc{Keywords} --- \textit{#1}}\par}

% --- L'appareil critique ----------------------------------------------------

% Le numéro de feuillet, en marge. C'est l'ancre qui permet de régler un
% désaccord sur une formule en regardant *un* feuillet plutôt qu'une chemise
% de six cents. Le site s'en sert aussi pour tourner les pages du fac-similé
% quand on fait défiler la transcription.
\newcommand{\leaf}[1]{%
  \marginnote{\footnotesize\color{gray!70}\textsf{#1}}%
  \label{leaf:#1}\ignorespaces}

% Illisible : on ne devine pas. Un mot inventé qui se lit comme les autres est
% le pire résultat possible.
\newcommand{\ill}{\textcolor{red!60!black}{[\,\ldots\,]}}

% Lecture douteuse : le mot est donné, et signalé comme incertain.
\newcommand{\uncertain}[1]{\uline{#1}}

% Ajout de l'éditeur — une lettre manquante, un mot sous-entendu.
\newcommand{\add}[1]{\textcolor{blue!50!black}{[#1]}}

% Biffé par Grothendieck lui-même. Ce qu'il a rayé fait partie du document :
% une rature dit souvent où la pensée a changé de direction.
\newcommand{\struck}[1]{\sout{#1}}

% Note du transcripteur, sur l'état matériel du feuillet.
\newcommand{\note}[1]{\par\smallskip{\small\color{gray!60!black}\textit{[#1]}}\par\smallskip}

% Note marginale de Grothendieck, distincte du corps du texte.
\newcommand{\marginal}[1]{\par\smallskip\noindent\hspace{1em}%
  {\small\color{gray!60!black}#1}\par\smallskip}

% --- Titre ------------------------------------------------------------------

\AtBeginDocument{%
  \begin{center}
    {\large\bfseries Fonds Alexandre Grothendieck --- cote n\textsuperscript{o}~\grothFolder}\\[2pt]
    {\itshape\grothFoldertitle}\\[4pt]
    {\small feuillets \grothFirst--\grothLast\ (lot \grothBatch)}\\[2pt]
    {\footnotesize\color{gray!60!black}Transcription établie d'après le fac-similé numérisé
     par l'Université de Montpellier.\\
     Les lectures incertaines sont soulignées, les illisibles notées [\,\ldots\,],
     les ajouts entre crochets.}
  \end{center}
  \vspace{1em}}

\endinput


\folder{161-1}
\batch{1}
\pages{1}{19}
\dating{s.d. — le groupe « Dossiers rassemblés par Grothendieck sur différentes thématiques » (161-1 à 162-6) est daté [après 1961-vers 1977]}
\foldertitle{Catégories : notes manuscrites (s.d.) — lecture modernisée}

\begin{document}

\section*{Résumé}

\begin{resume}
Une chemise étiquetée « Catégories » au crayon, dix-neuf pages sans
date, et quatre chantiers ouverts sur la même table : ces notes sont
l'atelier de Grothendieck au travail sur l'outillage général des catégories
— non pas un texte suivi, mais l'endroit où l'on vérifie que les outils
coupent.

Le premier chantier regarde ce que deux catégories se disent à travers un
couple de foncteurs adjoints — un aller et un retour, comme deux
dictionnaires entre deux langues. Traduire puis retraduire ne rend presque
jamais la phrase de départ ; mais les phrases qui reviennent inchangées
forment, de chaque côté, un noyau où les deux langues se répondent
exactement. Les pages 3 à 5 isolent ces deux noyaux, montrent que
l'aller-retour y devient une correspondance parfaite, et démontent
entièrement le cas où traduire deux fois ne perd plus rien après le premier
passage — ce qu'on appelle aujourd'hui une adjonction idempotente : elle
équivaut à la donnée de deux sous-catégories privilégiées et d'une
équivalence entre elles. Les pages 6 et 7 poursuivent avec une question
de reconnaissance : peut-on retrouver la sous-catégorie privilégiée comme
l'ensemble des objets « aveugles » aux flèches que la traduction écrase ?
« Or c'est vrai sauf erreur… », conclut la page — et c'est vrai.

Le deuxième chantier (pages 9 à 12) attaque une gêne familière :
multiplier un paquet fini d'objets sans choisir l'ordre des facteurs. Pour
des nombres la question est vide ; pour des objets, chaque échange de deux
facteurs laisse une trace, et la construction centrale — la catégorie
$\Phi(A)$ des familles finies, qu'on appelle aujourd'hui la catégorie
monoïdale symétrique libre — est précisément la machine qui rend ces traces
comptables. Grothendieck y ajoute des « contractions » (apparier deux
facteurs et les annihiler, comme un indice haut contre un indice bas en
calcul tensoriel), puis mesure l'obstruction exacte à un produit vraiment
sans ordre : un signe $\varepsilon(L)$ attaché à chaque objet inversible,
qui vaut $\pm 1$ et disparaît exactement dans ce qu'il nomme les catégories
de Picard.

Le troisième chantier (page 13, page 2 d'une note dont la première
manque) prend les théories elles-mêmes pour objets : une théorie n'est pas
une liste d'axiomes mais une catégorie, ses modèles sont des foncteurs, et
les « constructions » qu'on peut faire dans toute théorie forment à leur
tour une catégorie. La page est barré d'un grand trait oblique — le
geste de qui a reporté la note ailleurs.

Le dernier chantier (pages 15, 17, 19) est le plus proche des topos :
si l'on connaît un foncteur sur une petite partie « dense » d'une catégorie
— dense comme les rationnels dans les réels : tout objet est limite
d'objets de la partie — que faut-il pour le prolonger à tout en respectant
les limites ? La page 15 réfute d'abord une fausse piste par un
contre-exemple de deux lignes, puis retourne la question et touche un
théorème de reconnaissance : si le prolongement est exact, la grande
catégorie est un topos. C'est, à une page près, le critère de Giraud.

Rien n'est daté ; le ton est celui d'un carnet de questions — « demander
Schanuel », « Est-il vrai que… ? », « Réciproque ? » — où chaque réponse
ouvre la suivante. Les noms modernes sous lesquels chercher : adjonctions
idempotentes et sous-catégories réflexives, localisation et objets locaux,
catégorie monoïdale symétrique libre, catégories de Picard, théories de
Lawvere, extensions de Kan et théorème de Giraud.

\keywords{idempotent adjunction, reflective subcategory, free symmetric monoidal category, Picard category, Lawvere theory, left Kan extension}
\end{resume}

\section*{La partie fixe d'une adjonction (pages 3 à 5)}

Soit un couple de foncteurs adjoints
\[
E \underset{v}{\overset{u}{\rightleftarrows}} F,
\qquad \eta : \mathrm{id}_E \to vu, \quad \varepsilon : uv \to \mathrm{id}_F,
\]
$u$ adjoint à gauche de $v$.\footnote{Le manuscrit note $i$ l'unité et $j$
la coünité ; nous écrivons $\eta$ et $\varepsilon$, et gardons ses lettres
$u$, $v$.}

\subsection*{Restriction aux sous-catégories pleines}

Pour une sous-catégorie pleine $E' \subset E$ et $X, Y \in E'$,
l'application composée
\[
\mathrm{Hom}_E(X, Y) \longrightarrow \mathrm{Hom}_F(uX, uY)
\simeq \mathrm{Hom}_E(X, vuY)
\]
n'est autre que $\mathrm{Hom}(X, \eta_Y)$. Il en résulte aussitôt : si
$\eta_Y$ est un isomorphisme (resp. un monomorphisme) pour tout
$Y \in E'$, alors $u|E'$ est pleinement fidèle (resp. fidèle). En notant
$\overline{E}'$ la clôture de $E'$ par isomorphie, on a l'équivalence
\[
\text{(1)} \quad u|E' \text{ pl. fidèle et } vu(E') \subset \overline{E}'
\iff \forall\, Y \in E',\ \eta_Y : Y \xrightarrow{\ \sim\ } vu(Y),
\]
et, dualement, pour $F' \subset F$ pleine :
\[
\text{(2)} \quad v|F' \text{ pl. fidèle et } uv(F') \subset \overline{F}'
\iff \forall\, P \in F',\ \varepsilon_P : uv(P) \xrightarrow{\ \sim\ } P.
\]
Le manuscrit se demande en marge si la seconde condition de (1) est
vraiment nécessaire — « donner un exemple montrant que cette condition
n'est pas surabondante, avec $E'$ réduit à un élément ». Elle l'est : la
pleine fidélité de $u|E'$ ne force pas $vu$ à ramener $E'$ dans
$\overline{E}'$.

Les sous-catégories strictement pleines $E'$ vérifiant (1) et $F'$
vérifiant (2) se correspondent bijectivement par image essentielle,
$\alpha(E') = \overline{u(E')}$ et $\beta(F') = \overline{v(F')}$, et pour
un couple $(E', F')$ se correspondant, les foncteurs induits
$u' : E' \to F'$ et $v' : F' \to E'$ sont des équivalences quasi-inverses.

\subsection*{Les points fixes}

Posons
\[
E_0 = \{\, Y \in E \mid \eta_Y \text{ iso} \,\}, \qquad
F_0 = \{\, P \in F \mid \varepsilon_P \text{ iso} \,\},
\]
sous-catégories strictement pleines : la plus grande $E'$ vérifiant (1), la
plus grande $F'$ vérifiant (2). Alors
\[
F_0 = \overline{u(E_0)}, \qquad E_0 = \overline{v(F_0)},
\]
et $u$, $v$ induisent des équivalences quasi-inverses
$E_0 \simeq F_0$.\footnote{C'est l'énoncé aujourd'hui standard sur la
\emph{partie fixe d'une adjonction} : toute adjonction se restreint en une
équivalence entre les objets où l'unité, resp. la coünité, est inversible.}

\subsection*{Adjonctions idempotentes}

La page 5 dresse la liste des conditions équivalentes suivantes — on
dit aujourd'hui que l'adjonction est \emph{idempotente} lorsqu'elles sont
remplies :

\begin{enumerate}
\item $\mathrm{Im}\, v \subset E_0$, et de même $\mathrm{Im}\, u \subset F_0$ ;
\item $\eta v : v \xrightarrow{\ \sim\ } vuv$, et de même
$\varepsilon u : uvu \xrightarrow{\ \sim\ } u$ ;\footnote{Fait aujourd'hui
classique : chacune des quatre conditions ($\eta v$, $v\varepsilon$,
$\varepsilon u$, $u\eta$ inversible) implique les trois autres, de sorte
que les paires du manuscrit sont redondantes — mais l'énoncer par paires
symétriques est plus parlant.}
\item $\mathrm{Im}\, v \subset E_0$ et l'inclusion $E_0 \subset E$ admet un
adjoint à gauche ; de même $\mathrm{Im}\, u \subset F_0$ et l'inclusion
$F_0 \subset F$ admet un adjoint à droite ;
\item $u$ se factorise en $E \xrightarrow{u_1} F_1 \xrightarrow{\beta} F$
avec $\beta$ pleinement fidèle et $u_1$ une localisation, et $v$ se
factorise dualement.\footnote{« Passage à catégorie des fractions » dans le
manuscrit — c'est la localisation $E \to E[\Sigma^{-1}]$. La seconde ligne
de cette condition est barrée sur la page, avec en marge « à revoir
pour simplifier » ; nous la retenons parce qu'elle est correcte et que la page suivante s'en sert.}
\end{enumerate}

Sous ces conditions, l'inclusion $F_0 \subset F$ admet pour adjoint à
droite $u_0 \circ v'$ ($v' : F \to E_0$ la corestriction de $v$), et
l'inclusion $E_0 \subset E$ admet pour adjoint à gauche $v_0 \circ u'$
($u' : E \to F_0$ la corestriction de $u$) : $E_0$ est
\emph{réflexive} dans $E$, $F_0$ \emph{coréflexive} dans $F$. Les deux
petits calculs de Hom du manuscrit sont exactement les preuves modernes.
Et la structure complète se reconstruit : la donnée d'une adjonction
idempotente $(u, v)$ entre $E$ et $F$ équivaut à la donnée
\begin{itemize}
\item d'une sous-catégorie strictement pleine réflexive $E_0 \subset E$
(adjoint à gauche $\alpha'$ de l'inclusion $\alpha$),
\item d'une sous-catégorie strictement pleine coréflexive $F_0 \subset F$
(adjoint à droite $\beta'$ de l'inclusion $\beta$),
\item d'une équivalence $E_0 \underset{v_0}{\overset{u_0}{\rightleftarrows}} F_0$,
\end{itemize}
en posant $u = \beta u_0 \alpha'$ et $v = \alpha v_0 \beta'$.

Deux exemples dégénérés ouvrent la page 6 : si $E$ est ponctuelle, $v$
est l'unique foncteur $F \to E$, son adjoint à gauche existe si et
seulement si $F$ a un objet initial, et cet adjoint (qui pointe l'objet
initial) est pleinement fidèle ; dualement avec un objet final.

\section*{Localisations et objets locaux (pages 6 et 7)}

Soit $\alpha : E_1 \to E$ pleinement fidèle admettant un adjoint à gauche
$\alpha'$ — une sous-catégorie réflexive, autrement dit — de sorte que
$\alpha'$ est une localisation : $E_1 \simeq E[\Sigma^{-1}]$, où $\Sigma$
est la classe des flèches que $\alpha'$ inverse. Grothendieck pose alors
une question de descente d'adjoints : si $\beta : E_1 \to F$ est pleinement
fidèle et si $u = \beta\alpha'$ admet un adjoint à droite, $\beta$
en admet-il un ?

Il la réduit, en supposant $\mathrm{Ob}\,F = \mathrm{Ob}\,E_1 \cup \{a\}$,
à une question de représentabilité pure : \emph{un foncteur
$\psi : E_1^{\circ} \to \mathrm{Ens}$ tel que $\psi \circ \alpha'$ soit
représentable est-il représentable ?} Comme les foncteurs sur
$E_1^{\circ}$ s'identifient aux foncteurs sur $E^{\circ}$ qui inversent
$\Sigma$, cela devient une question de reconnaissance : si
$\xi \in E$ est tel que $\mathrm{Hom}_E(-, \xi)$ transforme les flèches de
$\Sigma$ en bijections — un objet \emph{$\Sigma$-local}, dirait-on
aujourd'hui — existe-t-il $\xi_1 \in E_1$ avec
\[
\mathrm{Hom}_E(X, \xi) \simeq \mathrm{Hom}_{E_1}(\alpha' X, \xi_1)
\quad (\simeq \mathrm{Hom}_E(X, \alpha\,\xi_1))
\]
fonctoriellement en $X$ — autrement dit, $\xi$ est-il dans l'image
essentielle de $\alpha$ ? La page vérifie que
\[
(f : X \to Y) \in \Sigma \iff
\mathrm{Hom}_E(Y, \alpha Z_1) \simeq \mathrm{Hom}_E(X, \alpha Z_1)
\ \text{ pour tout } Z_1 \in E_1,
\]
et conclut : « Or c'est vrai sauf erreur… ». C'est vrai, en effet, et c'est
aujourd'hui un lemme fondamental de la théorie des localisations
réflexives : l'image essentielle d'une sous-catégorie réflexive est
exactement la sous-catégorie des objets locaux pour les flèches inversées
par le réflecteur.\footnote{La réponse à la question de départ est alors
positive elle aussi. Le manuscrit note ce lieu des objets locaux
$\Phi(\Sigma)$ — notation sans rapport avec le $\Phi(A)$ des pages
suivants. Au passage, la page 7 écrit deux fois
$\mathrm{Hom}_{E_1}(Y, \alpha Z_1)$ là où seul $\mathrm{Hom}_E$ a un sens ;
lapsus de plume, corrigé ici.}

\section*{Le produit tensoriel sans parenthèses (pages 9 et 10)}

Le manuscrit écrit « $\otimes$-catégorie AUC » — associative, unitaire,
commutative — pour ce qu'on appelle aujourd'hui une \emph{catégorie
monoïdale symétrique}.\footnote{Le vocabulaire moderne et les théorèmes de
cohérence datent de 1963--1965 (Mac Lane, Bénabou) ; les pages ne sont
pas datés, et rien n'indique s'ils précèdent ou suivent ces publications.}

\subsection*{La catégorie des familles finies}

Pour une catégorie $A$, soit $\Phi(A)$ la catégorie des familles finies
d'objets de $A$ : un objet est un couple $(I, (L_i)_{i \in I})$, $I$
ensemble fini ; un morphisme
$(I, (L_i)) \to (J, (M_j))$ est un couple d'une bijection
$\tau : I \xrightarrow{\sim} J$ et d'isomorphismes
$\varphi_i : L_i \xrightarrow{\sim} M_{\tau(i)}$.\footnote{Le manuscrit ne
prend que des \emph{isomorphismes} $\varphi_i$ : son $\Phi(A)$ est donc le
groupoïde monoïdal symétrique libre sur le groupoïde des isomorphismes de
$A$. La variante aujourd'hui usuelle admet des $\varphi_i$ quelconques et
donne la catégorie monoïdale symétrique libre sur $A$ ; tout ce qui suit
vaut pour les deux, en lisant les Hom du bon côté.} La réunion disjointe
des familles,
\[
(I, (L_i)) \otimes (J, (M_j)) = (I \sqcup J,\ (L_i;\ M_j)),
\]
fait de $\Phi(A)$ une catégorie monoïdale symétrique : associativité par
celle de $\sqcup$, unité la famille vide, symétrie par l'échange des deux
cofacteurs — et tout cela sans utiliser la moindre structure sur $A$, comme
le note le manuscrit. Le cardinal découpe $\Phi(A)$ en tranches,
\[
\Phi(A) = \coprod_{n \geq 0} \Phi_n(A), \qquad
\Phi_0(A) = \{e\}, \quad \Phi_1(A) \simeq A.
\]

\textbf{Propriété universelle} (la Proposition de la page 9) : pour toute
catégorie monoïdale symétrique $C$, la restriction à $\Phi_1(A) \simeq A$
\[
\underline{\mathrm{Hom}}^{\otimes}(\Phi(A), C) \longrightarrow
\underline{\mathrm{Hom}}(A, C)
\]
est une équivalence de catégories — $\Phi(A)$ est \emph{libre}. La
démonstration est laissée en quatre points de suspension (« Dém. …. »).

\subsection*{Produits non parenthésés}

Le corollaire est le vrai but : pour toute catégorie monoïdale symétrique
$C$, il existe, à isomorphisme unique près, un unique foncteur monoïdal
symétrique
\[
\pi : \Phi(C) \longrightarrow C
\]
qui vaut l'identité sur $\Phi_1(C)$. Pour une famille finie
$(I, (L_i))$, l'objet $\pi((L_i))$ est \emph{le} produit tensoriel de la
famille — défini sans choisir d'ordre ni de parenthésage. Réciproquement,
la donnée d'un tel $\pi$, identité sur $\Phi_1(C)$ et rendant commutatif

\begin{tikzcd}[column sep=small, row sep=small]
\Phi_2(C) \arrow[r, "\pi\vert\Phi_2"] & C \\
C \times C \arrow[u] \arrow[ur, "\otimes"'] &
\end{tikzcd}

équivaut à la donnée d'une structure monoïdale symétrique sur $C$
prolongeant $\otimes$ : c'est la présentation \emph{non biaisée} des
catégories monoïdales, et l'équivalence des deux présentations est
précisément un théorème de cohérence.\footnote{Le manuscrit dit « il y a
une vérification » ; c'est la cohérence de Mac Lane sous sa forme non
biaisée. Noter aussi le NB de la page 10 : $\Phi$ est une monade —
$\pi$ appliqué dans $\Phi(C)$ lui-même aplatit une famille de familles en
la famille réunion disjointe.}

\section*{Contractions, inverses et le signe (pages 11 et 12)}

\subsection*{Familles à couples contractés}

La page 11 enrichit $\Phi(C)$ : soit $\Phi^{c}(C)$ la catégorie des
couples $(\xi, \gamma)$ où $\xi = (I, (L_i))$ est une famille finie et où
$\gamma$ consiste en un ensemble $\gamma'$ (éventuellement vide) de paires
disjointes $J = \{j', j''\} \subset I$, munies chacune d'un accouplement
parfait
\[
L_{j'} \otimes L_{j''} \longrightarrow \mathbf{1}_C
\]
— une dualité. Deux foncteurs d'oubli
$\mathrm{inc}, \mathrm{res} : \Phi^{c}(C) \to \Phi(C)$ gardent
respectivement toute la famille et la seule partie non contractée, et il
existe un morphisme canonique de foncteurs monoïdaux symétriques
\[
c\ :\ \pi \circ \mathrm{inc} \longrightarrow \pi \circ \mathrm{res},
\]
la \emph{contraction} : évaluer chaque paire appariée contre son
accouplement, comme on contracte un indice haut contre un indice bas en
calcul tensoriel. C'est, en substance, le calcul graphique des catégories
monoïdales à duaux — évaluations et « ficelles » comprises, la page 12
en dessine.

\subsection*{Objets inversibles et le signe $\varepsilon(L)$}

Pour deux objets $L_a$, $L_b$, les deux produits $L_a \otimes L_b$ et
$L_b \otimes L_a$ se correspondent par la symétrie ; mais si
$L_a = L_b = L$, la symétrie devient un \emph{automorphisme} de
$L \otimes L$, d'ordre 2, et il n'y a pas d'isomorphisme canonique entre
$L \otimes L$ et le produit non ordonné $\pi(L, L)$ de la famille à deux
éléments : en choisir un, c'est choisir un ordre total sur l'ensemble
d'indices.\footnote{Pour $n$ facteurs égaux il y a a priori $n!$
identifications de $\pi((L_i))$ avec $L^{\otimes n}$ ; la page écrit
« $L^{\otimes 2}$ » où l'on attend l'exposant $n$.}

Si $L$ est \emph{inversible}, l'obstruction se mesure : la symétrie de
$L \otimes L$ s'identifie à un élément
\[
\varepsilon(L) \in \mathrm{Aut}(\mathbf{1}_C), \qquad \varepsilon(L)^2 = 1,
\]
et le passage d'une identification $\pi(L,\dots,L) \simeq L^{\otimes n}$ à
une autre, via une permutation $\sigma$, se fait par
$\varepsilon(L)^{\mathrm{sign}(\sigma)}$. L'ambiguïté disparaît exactement
quand tous les $\varepsilon(L)$ valent $1$ : les \emph{catégories de
Picard} du manuscrit — tout objet inversible, tous les signes
triviaux.\footnote{Ce signe est devenu un invariant central : c'est lui qui
distingue les catégories de Picard \emph{strictes} et qui gouverne, par
exemple, la règle de Koszul des complexes gradués, où
$\varepsilon(L) = -1$ pour un inversible en degré impair.}

Enfin, l'\emph{inverse} d'un objet inversible $L$ : un $L'$ muni d'un
isomorphisme $L \otimes L' \simeq \mathbf{1}$ — d'où, par la symétrie, un
isomorphisme $L' \otimes L \simeq \mathbf{1}$ ; la relation est symétrique
et le couple $(L, L')$ peut se donner sans ordre, par un isomorphisme
$\pi(L, L') \simeq \mathbf{1}$ sur le produit non ordonné. La page 12
compare alors deux trivialisations de $L' \otimes L$ : celle obtenue de
$\varphi_{L,L'} : L \otimes L' \simeq \mathbf{1}$ en composant avec la
symétrie, $\varphi_{L',L} = \varphi_{L,L'} \circ s_{L',L}$, et celle,
$\psi_{L',L}$, définie par l'identité triangulaire
\[
\varphi_{L,L'} \otimes \mathrm{id}_L = \mathrm{id}_L \otimes \psi_{L',L}
\ :\ L \otimes L' \otimes L \longrightarrow L.
\]
« Relations entre $\varphi_{L',L}$ et $\psi_{L',L}$ ? » demande la page, qui s'interrompt sur le carré de naturalité de la symétrie et
trois colonnes de ficelles. La réponse — qu'il a visiblement en main —
est que les deux diffèrent précisément du facteur $\varepsilon(L)$ :
c'est la version « inverses » du signe du paragraphe précédent.

\section*{Théories et leurs constructions (page 13)}

Une page isolée — la page 2 d'une note dont la première page n'est pas
dans la chemise, barré d'un grand trait oblique — travaille sur des
\emph{théories} au sens catégorique : une théorie $T$ est une catégorie à
limites finies (resp. à limites quelconques), un modèle de $T$ dans une
catégorie $C$ à limites finies est un foncteur $T \to C$ exact à gauche, et
$T(C)$ désigne la catégorie des modèles.\footnote{C'est le point de vue des
\emph{théories de Lawvere} (1963) et de leurs généralisations à limites
finies — théories essentiellement algébriques, esquisses de Ehresmann. La page manipule ces objets comme allant de soi, en des termes qui lui
sont propres.} Les théories forment une 2-catégorie $\mathcal{T}$ (resp.
$\mathcal{T}_0$), et, un ensemble $I$ étant fixé, les $I$-morphismes de
théories en forment une autre, $\mathcal{T}^{(I)}$, avec un carré évident
de 2-foncteurs les comparant.

La notion mise en avant est celle de \emph{construction} pour une théorie
$T$ : un morphisme de théories
\[
T \longrightarrow \dot{\imath},
\]
où $\dot{\imath}$ est la théorie de l'« ensemble de base sans structure »,
celle dont les modèles dans $C$ redonnent $C$ : $\dot{\imath}_C = C$. Les
constructions forment une catégorie
\[
\Gamma^{T} = \underline{\mathrm{Hom}}(T, \dot{\imath}),
\]
qui a des limites finies (resp. des limites), munie d'une évaluation
$\Gamma^{T} \times T \to \dot{\imath}$ : pour tout modèle $\xi \in T(C)$,
le foncteur $w \mapsto w(\xi)$, de $\Gamma^{T}$ vers $C$, commute aux
limites en question. En langage d'aujourd'hui : $\Gamma^{T}$ est la
catégorie des \emph{opérations naturelles} sur les modèles de $T$ — ce que
toute structure de type $T$ sait fabriquer, uniformément en la structure —
et l'évaluation est le premier pas vers l'adjonction
structure--sémantique.\footnote{Le membre de droite du premier diagramme de la page, rogné par la marge, n'a pas pu être lu en entier ; la
transcription porte le détail.}

\section*{Prolongements : préfaisceaux, densité, topos (pages 15, 17, 19)}

\subsection*{Une fausse piste réfutée (page 15)}

Première question : $C$ petite à limites finies,
$\varphi : C \to \mathrm{Ens}$ exact à droite ; le prolongement
$\overline{\varphi}(\psi) = \mathrm{Hom}(\psi, \varphi)$ aux préfaisceaux
est-il encore exact à droite — transforme-t-il les limites projectives
finies de préfaisceaux en limites inductives d'ensembles ? Il faudrait pour
cela, note la page, que $\varphi$ soit un objet \emph{injectif} : que
tout morphisme $\psi \to \varphi$ s'étende le long de tout monomorphisme
$\psi \to \psi'$. En prenant $\psi = \varphi$ et
$\psi' = \varphi \sqcup \rho$, cela forcerait
$\mathrm{Hom}(\rho, \varphi) \neq \emptyset$ pour tout $\rho$. « Mais c'est
\emph{toujours} faux » : prendre pour $\rho$ un foncteur constant de valeur
non vide et un $\varphi$ à valeur vide quelque part —
$\mathrm{Hom}(x, \emptyset) = \emptyset$. Deux lignes, et la question est
close.

\subsection*{Le prolongement de Kan et la reconnaissance des topos}

Seconde question, celle qui porte : $C$ petite à limites finies, $E$ à
limites finies et petites colimites, $\varphi : C \to E$ exact à gauche,
et $\overline{\varphi} : \widehat{C} \to E$ le prolongement commutant aux
petites colimites — l'extension de Kan à gauche de $\varphi$ le long du
plongement de Yoneda. Sous quelles conditions sur $E$ le prolongement
$\overline{\varphi}$ est-il encore exact à gauche ? « On sait que c'est OK
si $E$ est un topos.\footnote{C'est le théorème aujourd'hui attaché au nom
de Diaconescu : sur un topos, l'extension de Kan d'un foncteur plat — en
particulier exact à gauche sur un $C$ à limites finies — est exacte à
gauche. La phrase de la page (« on sait que ») indique un fait déjà acquis
dans son entourage ; les pages ne sont pas datés.} Réciproque ? »

Et la réciproque est aussitôt esquissée : supposons $\varphi$ strictement
générateur — l'image de $C$ engendre $E$ par colimites, de façon dense.
Alors le foncteur $E \to \widehat{C}$, $X \mapsto \mathrm{Hom}(\varphi(-), X)$,
est pleinement fidèle et $\overline{\varphi}$ en est l'adjoint à gauche :
$E$ est une localisation réflexive de $\widehat{C}$. Donc si
$\overline{\varphi}$ est exact à gauche, $E$ est une localisation exacte à
gauche d'une catégorie de préfaisceaux — « $E$ est un topos ». C'est,
énoncé sur une demi-page, le théorème de reconnaissance de
Giraud.\footnote{Aux hypothèses de taille près, que la page ne discute
pas : il faut $C$ petite et $E$ à colimites petites, ce qui est supposé
ici. Une ligne intermédiaire, « il suffit que $E$ soit un topos », est
barrée sur la page — c'est bien la réciproque qu'il visait.}

\subsection*{Densité : le calcul complet (pages 17 et 19)}

Les deux derniers pages font le calcul général qui sous-tend tout cela.
Soient $C \subset E$ une sous-catégorie pleine strictement génératrice
— \emph{dense}, dit-on aujourd'hui : tout $X \in E$ est colimite canonique
$X = \varinjlim_{Y \in C/X} Y$ — et $F$ une catégorie à petites colimites.
L'extension de Kan et la restriction
\[
\underline{\mathrm{Hom}}(C, F)
\underset{\beta}{\overset{\alpha}{\rightleftarrows}}
\underline{\mathrm{Hom}}(E, F),
\qquad
\alpha(\varphi)(X) = \varinjlim_{Y \in C/X} \varphi(Y), \quad
\beta(\psi) = \psi|C,
\]
forment une adjonction, $\mathrm{Hom}(\alpha\varphi, \psi) \simeq
\mathrm{Hom}(\varphi, \beta\psi)$, avec $\beta\alpha \simeq \mathrm{id}$ :
$\alpha = i_!$ est pleinement fidèle, comme $i$.\footnote{L'unité
$\beta\alpha \simeq \mathrm{id}$ utilise seulement que $C$ est pleine dans
$E$ ; la densité stricte sert plus bas, pour
$\alpha\beta(\psi)(X) = \varinjlim_{C/X} \psi(Y) \simeq \psi(X)$.}

Quand $\psi$ commute aux petites colimites, on a
$\alpha\beta(\psi) \simeq \psi$ puisque $X = \varinjlim_{C/X} Y$. Notant
$\underline{\mathrm{Hom}}'(E, F)$ les foncteurs commutant aux petites
colimites, et $\underline{\mathrm{Hom}}'(C, F)$ les $\varphi$ vérifiant la
condition de compatibilité

(*) \emph{pour tout système inductif $(Y_i)$ de $C$ dont la colimite dans
$E$ appartient encore à $C$, $\varphi$ envoie cette colimite sur
$\varinjlim \varphi(Y_i)$,}

la restriction induit un foncteur pleinement fidèle
\[
\beta' : \underline{\mathrm{Hom}}'(E, F) \longrightarrow
\underline{\mathrm{Hom}}'(C, F),
\]
et la question est de savoir si c'est une équivalence — c'est-à-dire si
$\alpha$ envoie $\underline{\mathrm{Hom}}'(C, F)$ dans
$\underline{\mathrm{Hom}}'(E, F)$. Par un système inductif double, tout se
ramène au

\textbf{Lemme.} \emph{Si $\varphi$ vérifie (*), alors pour tout
système inductif $(Y_i)$ dans $C$ ayant une colimite $Y$ dans $E$, le
morphisme canonique $\varinjlim \varphi(Y_i) \to \alpha\varphi(Y)$ est un
isomorphisme.}\footnote{La page écrit $\beta\varphi(Y)$ — lapsus pour
$\alpha\varphi(Y)$, seul défini quand $Y$ n'est pas dans $C$ ; la
transcription le signale. C'est le lemme de comparaison des colimites le
long d'un foncteur dense, au cœur de ce qu'on appelle aujourd'hui le
théorème de densité.}

Cela revient à comparer les colimites de $\varphi$ sur les deux systèmes
$(Y_i)$ et $C/Y \to C$, donc à la \emph{cofinalité} du foncteur
$I \to C/Y$. « Cela est-il cofinal ? C'est vrai en tout cas si pour tout
diagramme fini dans $C$, la colimite dans $E$ existe et est dans $C$. » Un
point d'interrogation en marge ferme le dossier — la question de
cofinalité, en général, reste ouverte sur ces pages.

\end{document}
